\documentclass[11pt,a4paper]{article}
\usepackage[T1]{fontenc}
\usepackage[utf8]{inputenc}
\usepackage{lmodern}
\usepackage[margin=25.5mm,headheight=15pt]{geometry}
\usepackage{amsmath,amssymb,amsthm,mathtools}
\usepackage{booktabs,longtable,tabularx,array,microtype,enumitem,xcolor}
\usepackage{fancyhdr,listings,xurl,titlesec}
\definecolor{navy}{RGB}{30,57,82}
\definecolor{ink}{RGB}{32,38,44}
\definecolor{pale}{RGB}{239,243,247}
\usepackage[colorlinks=true,linkcolor=ink,citecolor=ink,urlcolor=navy,
  bookmarksnumbered=true]{hyperref}
\usepackage{bookmark}
\hypersetup{pdftitle={Perfect Exact Pairs in Coarse-Description Classes},
  pdfauthor={Research note prepared for Vladimir Reshetnikov},
  pdfsubject={Coarse reducibility, least Turing representatives, cone avoidance and exact pairs}}
\urlstyle{same}
\titleformat{\section}{\large\bfseries\color{navy}}{\thesection}{0.7em}{}
\titleformat{\subsection}{\normalsize\bfseries\color{navy}}{\thesubsection}{0.65em}{}
\setlength{\parskip}{3.5pt}
\setlength{\parindent}{1.2em}
\setlength{\emergencystretch}{3em}
\setlist{nosep,leftmargin=*}
\setcounter{tocdepth}{1}
\pagestyle{fancy}\fancyhf{}
\fancyhead[L]{\small\color{navy}Exact pairs and coarse descriptions}
\fancyhead[R]{\small September 2026}
\fancyfoot[C]{\thepage}
\renewcommand{\headrulewidth}{0.3pt}
\newcommand{\N}{\mathbb N}
\newcommand{\Cantor}{2^{\N}}
\newcommand{\Baire}{\N^{\N}}
\newcommand{\leT}{\leq_{\mathrm T}}
\newcommand{\nleT}{\nleq_{\mathrm T}}
\newcommand{\eqT}{\equiv_{\mathrm T}}
\newcommand{\degT}{\operatorname{deg}_{\mathrm T}}
\newcommand{\nc}{\mathrm{nc}}
\newcommand{\uc}{\mathrm{uc}}
\newcommand{\CD}{\mathcal C}
\newcommand{\core}{\mathcal I}
\newcommand{\Comp}{\mathsf{Comp}}
\newcommand{\low}[1]{\mathord\downarrow_{\!\mathrm T}#1}
\newcommand{\AP}{\mathsf{AP}}
\newcommand{\patch}{\mathbin\star}
\newcommand{\cl}[1]{\overline{#1}}
\newcommand{\Spec}{\operatorname{Spec}_{\mathrm T}}
\newcommand{\restr}{\mathbin\upharpoonright}
\newcommand{\pair}[2]{\langle #1,#2\rangle}
\DeclareMathOperator{\dom}{dom}
\DeclareMathOperator{\diam}{diam}
\newtheorem{theorem}{Theorem}[section]
\newtheorem{lemma}[theorem]{Lemma}
\newtheorem{proposition}[theorem]{Proposition}
\newtheorem{corollary}[theorem]{Corollary}
\theoremstyle{definition}\newtheorem{definition}[theorem]{Definition}
\theoremstyle{remark}\newtheorem{remark}[theorem]{Remark}
\lstset{basicstyle=\small\ttfamily,columns=fullflexible,breaklines=true,
  frame=single,showstringspaces=false}
\begin{document}
\thispagestyle{empty}
\noindent{\small\sffamily\bfseries\color{navy}COMPUTABILITY THEORY / RESEARCH NOTE}\par
\vspace{7mm}
{\setlength{\parindent}{0pt}\fontsize{25}{29}\selectfont\bfseries\color{navy}Perfect Exact Pairs in\par Coarse-Description Classes\par}
\vspace{4mm}
\noindent{\large An attack on the least-Turing-representative question}\par
\vspace{5mm}
\noindent Prepared for Vladimir Reshetnikov\\[2pt]
18 September 2026\par
\vspace{5mm}
\begin{abstract}
We investigate the least-representative question selected as C1 in the supplied
Turing-degree survey. For a set $X$, let $\CD_X$ be the sets differing from $X$
on a density-zero set, and let $\core_X$ be the sets computable from every
member of $\CD_X$. We give a conventional proof that $\CD_X$ contains a perfect
family of pairwise Turing-incomparable members such that the common lower cone
of any two is exactly $\core_X$. The construction can simultaneously control
intersections with the lower cones of any prescribed countable family of
oracles. Its main ingredient is a local recovery lemma in a complete metric
on $\CD_X$: information computed on a dense open-region subset is already
computable from every coarse description.

An explicit computably enumerable, many-one complete set $X$ has
$\core_X=\Comp$ but no computable coarse description. Consequently its uniform
and nonuniform coarse-equivalence classes contain Turing minimal pairs, have
no least Turing degree, and have no countable coinitial family of Turing
degrees. Relativized examples have any prescribed principal core. All infinite
arguments are written out; the accompanying program checks only finite
arithmetic lemmas.
\end{abstract}
\vspace{3mm}
\noindent\colorbox{pale}{\parbox{\dimexpr\linewidth-2\fboxsep\relax}{\small
\textbf{Status and attribution.} The negative answer to the literal C1 question
already follows from results published in 2016; it is not claimed here as a
new resolution of an open problem. The perfect-family and simultaneous
exact-pair development is the stronger result pursued in this note. Its proof
is supplied in full, but its novelty has not been established, and it has not
been independently refereed or checked by a proof-assistant kernel.}}
\vfill
{\noindent\small Conventional proofs, an explicit c.e. witness, and reproducible finite checks.\par}
\normalsize
\newpage
{\small\tableofcontents}
\newpage

\section{Target, result, and correction of the research map}
\label{sec:target}
The uploaded report~\cite{Survey}, entry C1, asks whether every nonuniform
coarse-equivalence class contains a representative of least Turing degree.
For total functions, its statement is
\begin{equation}\label{eq:question}
 \forall f\in\Baire\ \exists g\equiv_{\nc}f\
 \forall h\equiv_{\nc}f\quad g\leT h.
\end{equation}
The report proposes a concrete negative route: find a nonzero coarse class
containing two representatives that form a Turing minimal pair. This is the
route taken here, with an additional perfect-set strengthening.

\subsection{What the source check changes}
Gerdes's 2025 preprint explicitly asks the least-degree question in
Question~7~\cite{Gerdes}. However, the invariant used below is already
Definition~3.1 of Hirschfeldt--Jockusch--Kuyper--Schupp~\cite{HJKS}, where it is
written $X^{\mathrm c}$. Their Theorem~4.3, attributed to Igusa, gives a trivial
invariant when the coarse computability bound is $1$, including examples that
are not coarsely computable. Those two properties alone negate
\eqref{eq:question}, as the next paragraph proves. The journal publication is
from 2016. Thus the literal C1 entry is not a defensible novelty target,
despite its appearance in the later question list.

Suppose every coarse description of $X$ has only computable information in
common, while $X$ has no computable coarse description. Any least
representative $g$ of the coarse class would be computable from every coarse
description of $X$, since each such description itself belongs to the class.
The graph of $g$ would therefore be computable. But a computable member of the
class would provide a computable coarse description of $X$, a contradiction.
This argument works for both uniform and nonuniform coarse equivalence.

The status correction does not end the investigation. The older invariant
statement does not itself present the minimal-pair certificate requested by
the survey, nor the exact-pair and perfect-family results proved here. The
remainder develops those stronger statements without relying on randomness,
a priority minimal-pair theorem, or an imported cone-avoidance theorem.

\subsection{The principal structural theorem}
Write $\low B=\{A\subseteq\N:A\leT B\}$. The word \emph{exact} refers to equality
of common lower cones with an entire ideal, not merely to the exclusion of
one selected noncomputable set.

\begin{theorem}[Perfect exact pairs; proved in Section~\ref{sec:perfect}]
\label{thm:mainpreview}
For every $X\subseteq\N$, put
\[
 \CD_X=\{D\subseteq\N:\rho_n(D\triangle X)\longrightarrow0\},
 \qquad \core_X=\bigcap_{D\in\CD_X}\low D.
\]
There is a perfect set $P\subseteq\CD_X$ of pairwise Turing-incomparable
members such that
\begin{equation}\label{eq:perfect-main}
 D,E\in P,\ D\ne E
 \quad\Longrightarrow\quad
 (\low D)\cap(\low E)=\core_X.
\end{equation}
Moreover, given any countable family of oracles $(B_j)_{j\in\N}$, $P$ may be
chosen so that, for every $D\in P$ and every $j$,
\begin{equation}\label{eq:prescribed-main}
 (\low D)\cap(\low {B_j})=\core_X\cap(\low {B_j}).
\end{equation}
No effective presentation of $(B_j)$ is assumed or needed.
\end{theorem}

Taking $B_0=X$ gives an exact partner for the original, prescribed $X$, not
merely two newly chosen descriptions. For the explicit $X$ constructed in
Section~\ref{sec:diagonal}, $X\eqT\emptyset'$ and $\core_X=\Comp$. Thus every
member of the selected perfect family forms a Turing minimal pair with
$\emptyset'$, as well as with every other member of that family.

\subsection{What is and is not being claimed}
The negative answer to C1 is a consequence of established work, not a new
breakthrough. The mathematical contribution attempted here is the fully
proved perfect exact-pair formulation and its applications. A targeted
literature search did not establish whether these stronger formulations have
already appeared, perhaps as consequences of generic-description arguments.
Their novelty therefore remains unverified.

The effective-dense problem C2 is not answered. Nor do the results prove that
the exhibited coarse classes contain no \emph{minimal} Turing degrees. A least
degree is below every representative; a minimal degree need only have no
strictly smaller representative. These are different assertions.

\section{Descriptions, reducibility, and the common-information ideal}
\label{sec:definitions}
Identify a set with its characteristic function. For $n\ge1$, define
\[
 \rho_n(S)=\frac{|S\cap[0,n)|}{n},\qquad
 \delta(Y,Z)=\limsup_{n\to\infty}\rho_n(Y\triangle Z).
\]
Thus $D\in\CD_X$ exactly when $\delta(D,X)=0$. Because prefix discrepancies
are nonnegative, a zero limsup is equivalent to convergence to zero.

For total functions $f,g$, define $f\le_{\nc}g$ to mean that every coarse
description of $g$ computes some coarse description of $f$. In the uniform
version $f\le_{\uc}g$, one Turing functional must work for every source
description. These are the conventions of the supplied report and
Gerdes~\cite{Gerdes}. We write $[X]_r$ for the equivalence class under $r$, with
$r\in\{\nc,\uc\}$.

Binary source sets cause no ambiguity between binary and natural-number-valued
descriptions. Replace every nonbinary answer in a description of $X$ by zero.
This transformation is computable and does not create a disagreement at any
previously correct position. Consequently the binary descriptions suffice for
the common-information questions below. Total functions are treated as oracles
by coding their graphs as sets; this coding preserves Turing degree.

\begin{lemma}[Equality of description classes]\label{lem:sameclass}
If $D\in\CD_X$, then $\CD_D=\CD_X$. In particular
$D\equiv_{\uc}X$, and hence $D\equiv_{\nc}X$.
\end{lemma}
\begin{proof}
For each prefix, disagreement with $X$ is contained in the union of
disagreement with $D$ and $D\triangle X$. Taking upper densities gives
\[
 \delta(E,X)\le\delta(E,D)+\delta(D,X).
\]
The same inequality with $X,D$ exchanged proves the equivalence of the two
zero-discrepancy conditions. The identity functional is a coarse reduction
in both directions. The argument also applies to natural-number-valued
coarse descriptions.
\end{proof}

\begin{definition}[Common-information ideal]
For $X\subseteq\N$, define
\[
 \core_X=\{A\subseteq\N:(\forall D\in\CD_X)\ A\leT D\}.
\]
We call this the \emph{core} only as an abbreviation for the displayed
invariant; the invariant itself is established terminology and notation
$X^{\mathrm c}$ in~\cite[Definition~3.1]{HJKS}.
\end{definition}

\begin{lemma}\label{lem:ideal}
The family $\core_X$ is a countable Turing ideal: it contains every computable
set, is downward closed under $\leT$, and is closed under finite Turing joins.
It is contained in $\low X$. If $X\le_{\nc}Y$, then
$\core_X\subseteq\core_Y$. In particular the core is invariant under both
coarse equivalences.
\end{lemma}
\begin{proof}
Every oracle computes every computable set. If $A\leT B$ and every description
computes $B$, every description computes $A$. If every description computes
both $A$ and $B$, it computes their join. Since $X$ itself is a description of
$X$, each member of the core is computable from $X$. There are only countably
many oracle programs, so $\low X$ and hence $\core_X$ are countable.

For monotonicity, fix $A\in\core_X$ and $D\in\CD_Y$. The assumed reduction
provides an $E\in\CD_X$ with $E\leT D$. Then $A\leT E\leT D$. Since this holds
for every $D$, $A\in\core_Y$. Applying this in both directions proves
invariance. Uniform reducibility implies nonuniform reducibility.
\end{proof}

\begin{remark}[Two different equivalence classes]
The space $\CD_X$ is the equivalence class of $X$ modulo density-zero
symmetric difference. It is a subset, generally a proper subset, of the
coarse-reducibility class $[X]_r$. We construct witnesses in the smaller
space. Confusing these two classes would invalidate the least-representative
argument.
\end{remark}

\section{A complete metric on one coarse-description class}
\label{sec:metric}
The asymptotic discrepancy $\delta$ is identically zero on $\CD_X\times\CD_X$,
so it cannot provide a useful topology there. Instead use every prefix:
\begin{equation}\label{eq:dmetric}
 d(Y,Z)=\sup_{n\ge1}\rho_n(Y\triangle Z).
\end{equation}
This is a metric on all of $\Cantor$. We use its restriction to $\CD_X$;
this restriction is not the topology inherited from ordinary Cantor space.

For a finite binary string $\sigma$ and a set $Z$, write $\sigma\patch Z$ for
$Z$ with its first $|\sigma|$ bits replaced by $\sigma$.

\begin{lemma}[Finite-prefix facts]\label{lem:prefix}
The following statements hold.
\begin{enumerate}[label=(\roman*)]
\item If $Y(j)\ne Z(j)$, then $d(Y,Z)\ge1/(j+1)$. Hence each ordinary finite
cylinder is open in the $d$ topology.
\item If $L=|\sigma|>0$, then
\begin{equation}\label{eq:patchidentity}
 d(\sigma\patch Z,Z)
 =\max_{1\le n\le L}\frac{|\{j<n:\sigma(j)\ne Z(j)\}|}{n}.
\end{equation}
For the empty string this distance is zero.
\item The triangle inequality holds for $d$ and for $\delta$.
\end{enumerate}
\end{lemma}
\begin{proof}
A disagreement at $j$ contributes at least one error to the prefix of length
$j+1$. In particular a $d$-ball of radius less than $1/L$ fixes the first $L$
bits. For (ii), after position $L$ the number of disagreements is constant,
so its proportion cannot exceed the proportion at $L$. For (iii), use the
inclusion
$Y\triangle W\subseteq(Y\triangle Z)\cup(Z\triangle W)$ at each prefix,
and take a supremum or a limsup. Symmetry is immediate, and (i) supplies
separation, completing the metric verification.
\end{proof}

\begin{proposition}[Polish description space]\label{prop:polish}
The metric space $(\CD_X,d)$ is nonempty, complete, separable, and has no
isolated points. For every $Z\in\CD_X$, the finite modifications of $Z$ are
dense in $(\CD_X,d)$.
\end{proposition}
\begin{proof}
Nonemptiness follows from $X\in\CD_X$. Let $(D_k)$ be $d$-Cauchy. For each
coordinate $j$, the first part of Lemma~\ref{lem:prefix} shows that $D_k(j)$
is eventually constant. Let $D(j)$ be this eventual value.

Given $\varepsilon>0$, choose $K$ so that $d(D_k,D_l)<\varepsilon/2$ whenever
$k,l\ge K$. Fix $k\ge K$ and a prefix length $n$. For sufficiently large $l$,
$D_l\restr n=D\restr n$, whence
$\rho_n(D_k\triangle D)\le\varepsilon/2$. This holds for every $n$, so
$d(D_k,D)\le\varepsilon/2<\varepsilon$. Thus $D_k\to D$ in $d$. Also
\[
 \delta(D,X)\le d(D,D_k)+\delta(D_k,X)=d(D,D_k),
\]
which tends to zero. Therefore $D\in\CD_X$, proving completeness.

Fix $Y,Z\in\CD_X$. Put $Z_N=(Y\restr N)\patch Z$, a finite modification of
$Z$. The first $N$ bits of $Z_N$ and $Y$ agree, and for $n>N$ their errors are
a subset of $Y\triangle Z$ below $n$. Therefore
\begin{equation}\label{eq:finite-dense}
 d(Z_N,Y)\le\sup_{n>N}\rho_n(Y\triangle Z)\longrightarrow0.
\end{equation}
The convergence holds because $\delta(Y,Z)=0$. Thus the countable set of
finite modifications of $Z$ is dense, proving separability as well.

Finally, flip just bit $j$ of an arbitrary $Y\in\CD_X$. The result remains
in $\CD_X$ and has distance exactly $1/(j+1)$ from $Y$. Taking $j$ arbitrarily
large proves that no point is isolated.
\end{proof}

A set is \emph{nowhere dense} if its closure has empty interior, and
\emph{meager} if it is a countable union of nowhere dense sets. We use the
Baire category theorem in the form: a nonempty complete metric space cannot
be covered by countably many nowhere dense sets. Its usual nested-ball proof
also shows that the complement of such a union is dense. These topological
statements are used with $d$, not with $\delta$.

\section{Local recovery and simultaneous cone avoidance}
\label{sec:recovery}
Fix an enumeration $(\Phi_e)$ of oracle Turing functionals. For a set $A$,
write
\[
 S_{e,A}=\{D\in\CD_X:\Phi_e^D=A\}.
\]
Equality here includes totality on every input and binary output. A
computation $\Phi_e^\sigma(n)\downarrow=v$ uses only oracle bits within the
finite string $\sigma$.

\begin{lemma}[Local recovery]\label{lem:recovery}
If $S_{e,A}$ is dense in a nonempty $d$-open set $U\subseteq\CD_X$, then
$A\in\core_X$.
\end{lemma}
\begin{proof}
Fix any $Z\in\CD_X$. Finite modifications of $Z$ are dense, so choose a finite
modification $Z^*$ in $U$. Choose a positive rational $q$ with
$B_d(Z^*,q)\subseteq U$, where the ball is taken within $\CD_X$. The finite
modification and $q$ may depend on $Z$; they will be constants in the
reduction below.

To compute $A(n)$ using $Z^*$, dovetail over finite strings $\sigma$ and
finite simulations of $\Phi_e^\sigma(n)$. Accept a convergent simulation
with value $v$ only if
\begin{equation}\label{eq:searchcondition}
 \frac{|\{j<k:\sigma(j)\ne Z^*(j)\}|}{k}<q
 \qquad(1\le k\le|\sigma|).
\end{equation}
Each test is a finite, $Z^*$-decidable calculation. No test of membership in
$\CD_X$, totality, or asymptotic density is performed.

\emph{Termination.} Density gives a $D\in S_{e,A}\cap B_d(Z^*,q)$. A finite
initial segment $\sigma\prec D$ witnesses convergence on input $n$. Every
prefix discrepancy of $D$ from $Z^*$ is at most $d(D,Z^*)<q$, so this $\sigma$
passes \eqref{eq:searchcondition}. The dovetailing eventually finds an
accepted computation.

\emph{Correctness.} For any accepted $\sigma$, the set
$W=\sigma\patch Z^*$ belongs to $\CD_X$, because it is a finite modification
of $Z^*$. By \eqref{eq:patchidentity} and the finitely many strict inequalities
in \eqref{eq:searchcondition}, $d(W,Z^*)<q$. Thus $W\in U\cap[\sigma]$.
The intersection is a nonempty $d$-open set. Density of $S_{e,A}$ in $U$
provides $D'\in S_{e,A}\cap U\cap[\sigma]$. The same finite computation runs
with oracle $D'$, so the accepted value is $v=A(n)$.

The search computes $A$ from $Z^*$. Since $Z^*$ is a fixed finite modification
of $Z$, $Z^*\eqT Z$, and therefore $A\leT Z$. The choice of $Z$ was arbitrary,
which proves $A\in\core_X$.
\end{proof}

\begin{remark}[The nonuniformity is essential]
The proof asserts $\forall Z\in\CD_X\ \exists i\ (\Phi_i^Z=A)$. It does not
supply one index working for every $Z$, an effective choice of the finite
patch, or an effective choice of the rational radius. Hard-coding finitely
many constants proves an ordinary Turing reduction; finding those constants
uniformly would be a stronger assertion.
\end{remark}

\begin{theorem}[All-or-meager cone dichotomy]\label{thm:dichotomy}
For any $A\subseteq\N$, the cone
$\{D\in\CD_X:A\leT D\}$ is either all of $\CD_X$ or a meager subset of it.
More precisely, if $A\notin\core_X$, each $S_{e,A}$ is nowhere dense.
\end{theorem}
\begin{proof}
When $A\in\core_X$, the cone is all of the space by definition. Otherwise,
if $\cl{S_{e,A}}$ had nonempty interior $U$, $S_{e,A}$ would be dense in $U$.
Lemma~\ref{lem:recovery} would imply $A\in\core_X$, a contradiction. Thus
each $S_{e,A}$ is nowhere dense. The cone is their countable union.
\end{proof}

\begin{corollary}[Countable avoidance]\label{cor:countavoid}
Given a countable family $(A_j)$ with every $A_j\notin\core_X$, the set
\[
 \{D\in\CD_X:(\forall j)\ A_j\nleT D\}
\]
is comeager, and in particular is nonempty and dense. No effective enumeration
of the family is required.
\end{corollary}
\begin{proof}
Remove all sets $S_{e,A_j}$. This removes a countable union of nowhere dense
sets. Proposition~\ref{prop:polish} and the Baire category theorem apply.
\end{proof}

\begin{theorem}[Exact partners for prescribed oracles]\label{thm:partners}
For every countable family $(B_j)$ there is $D\in\CD_X$ such that
\[
 (\low D)\cap(\low {B_j})=\core_X\cap(\low {B_j})
 \qquad\text{for all }j.
\]
The set of suitable $D$ is comeager. In particular, some $D\in\CD_X$ satisfies
\begin{equation}\label{eq:fixedexact}
 (\low D)\cap(\low X)=\core_X.
\end{equation}
\end{theorem}
\begin{proof}
There are countably many sets computable from any one oracle $B_j$, and hence
countably many in
\[
 \mathcal A=\bigcup_j\bigl((\low {B_j})\setminus\core_X\bigr).
\]
Avoid all members of $\mathcal A$ by Corollary~\ref{cor:countavoid}. If $A$ is
computable from both the resulting $D$ and $B_j$, it cannot belong to
$(\low {B_j})\setminus\core_X$, and thus belongs to the right-hand side.
Conversely every member of $\core_X$ is computable from $D$, which proves the
reverse inclusion. Take $B_0=X$ and use $\core_X\subseteq\low X$ for
\eqref{eq:fixedexact}.
\end{proof}

This avoids an important logical shortcut. The assertion that each individual
non-core set can be avoided does not, by itself, give one oracle avoiding all
of them simultaneously. The nowhere-density proof and Baire category are the
additional ingredients that make simultaneous avoidance valid.

\section{A perfect family of exact pairs}
\label{sec:perfect}
We now prove Theorem~\ref{thm:mainpreview}, including the word \emph{perfect}.
We give the finite-level fusion directly, rather than invoking a perfect-set
selection theorem as an unexplained black box.

\begin{lemma}[Bad pair relations are nowhere dense]\label{lem:badpairs}
For $e,f\in\N$, put
\[
 R_{e,f}=\{(D,E)\in\CD_X^2:
   \Phi_e^D=\Phi_f^E=A\in\Cantor\text{ for some }A\notin\core_X\}.
\]
Then $R_{e,f}$ is nowhere dense in the product of the $d$ topologies.
\end{lemma}
\begin{proof}
Suppose to the contrary that $R_{e,f}$ is dense in a nonempty open rectangle
$U\times V$. Such a rectangle would exist if its closure had interior.
Choose $(D_0,E_0)\in R_{e,f}\cap(U\times V)$, and let
$A=\Phi_e^{D_0}=\Phi_f^{E_0}\notin\core_X$.

Fix $n$. We claim that every convergent computation $\Phi_e^D(n)$ with $D\in U$
has value $A(n)$. Otherwise there would be a $D\in U$ whose computation
converges to a different value $v$. Choose finite initial segments witnessing
this computation and $\Phi_f^{E_0}(n)=A(n)$. Intersect their cylinders with
$U$ and $V$. The resulting nonempty open rectangle consists of pairs with
unequal outputs on input $n$, and is therefore disjoint from $R_{e,f}$. This
contradicts the assumed density. The claim holds for every $n$.

The projection of $R_{e,f}\cap(U\times V)$ onto $U$ is dense in $U$: every
nonempty open subset of $U$, multiplied by $V$, meets $R_{e,f}$. Every oracle
in this projection makes $\Phi_e$ total. By the preceding claim its output is
exactly $A$. Therefore $S_{e,A}$ is dense in $U$.
Lemma~\ref{lem:recovery} gives $A\in\core_X$, contradicting its choice.
\end{proof}

\begin{theorem}[Perfect exact-pair theorem]\label{thm:perfect}
There exists a continuous injection
$F:2^{\N}\longrightarrow(\CD_X,d)$ whose image $P$ satisfies
\eqref{eq:perfect-main}. It can be chosen to satisfy
\eqref{eq:prescribed-main} for any prescribed countable family $(B_j)$, and
so that $P\cap\core_X=\varnothing$. Its image is perfect in the ordinary
Cantor topology, and its members are pairwise Turing incomparable.
\end{theorem}
\begin{proof}
For each $(e,f)$ let
\[
 O_{e,f}=\CD_X^2\setminus\cl{R_{e,f}}.
\]
By Lemma~\ref{lem:badpairs}, these are open dense sets. They are indexed by
\emph{ordered} pairs of functionals; there is no symmetry assumption.

We also need countably many unary open dense requirements. For each set
$A$ computable from some $B_j$ but not in $\core_X$, and each $e$, include
$\CD_X\setminus\cl{S_{e,A}}$. Include as well
$\CD_X\setminus\{A\}$ for every $A\in\core_X\cap\CD_X$. These last sets are
open dense because singletons are closed and no point is isolated. There
are only countably many requirements, by Lemma~\ref{lem:ideal}. Enumerate
them as $(V_i)$, repeating the whole space if necessary. This enumeration is
an ordinary set-theoretic enumeration, not a computable one.

Construct, for each finite binary string $s$, a nonempty open metric ball
$U_s\subseteq\CD_X$. Write $K_s=\cl{U_s}$, with closure within $\CD_X$.
Require
\begin{align}
 &K_{s0}\cup K_{s1}\subseteq U_s,
       \qquad K_{s0}\cap K_{s1}=\varnothing,\label{eq:fusion1}\\
 &\diam(K_s)\le2^{-|s|}\quad\text{if }|s|>0,\label{eq:fusion2}\\
 &K_s\subseteq V_i\quad\text{if }i\le m=|s|,\ m\ge1,\label{eq:fusion3}\\
 &K_s\times K_t\subseteq O_{e,f}\quad
   \text{if }s\ne t,\ |s|=|t|=m\ge1,\ e,f\le m.\label{eq:fusion4}
\end{align}
Start with any nonempty open ball for the root; impose no diameter condition
at level zero.

Here is the complete extension step. In each current parent ball choose two
distinct points. Small enough balls around them have disjoint closures
contained in the parent and meet the next diameter bound. Such points exist
because the space has no isolated points. This gives preliminary child balls
at the next level.

There are finitely many children and finitely many unary requirements with
$i\le m$ at level $m$. Intersect each child successively with those open dense
sets, choosing a smaller nonempty ball whose closure lies in the intersection.
Then list all finitely many ordered distinct child pairs $(s,t)$ and all
$e,f\le m$. For the next pair requirement, the product of the current two
balls is a nonempty open rectangle, so it meets $O_{e,f}$. Because $O_{e,f}$
is open, choose smaller balls in the two coordinates whose product of
closures is contained in $O_{e,f}$. Shrink just those two child balls. Every
previously established requirement persists under shrinking. After finitely
many steps all requirements at that level hold. This completes the recursive
construction of the balls.

For $p\in2^{\N}$, the sets $(K_{p\restr m})_{m\ge1}$ are nonempty closed,
nested, and have diameters tending to zero. To see that they have a unique
common point without any compactness assumption, select one point from each.
The nesting and diameter bounds make this sequence Cauchy. Completeness puts
its limit in every closed set, and the vanishing diameters give uniqueness.
Define $F(p)$ to be that point. Shared long prefixes make the images close,
so $F$ is continuous. Disjoint sibling closures make it injective.

Let $p\ne q$, and fix $(e,f)$. At a sufficiently large level $m$, the two
prefixes are distinct and $m\ge e,f$. Requirement~\eqref{eq:fusion4} gives
$(F(p),F(q))\in O_{e,f}$, and thus this pair is not in $R_{e,f}$. If a set $A$
is computable from both images, some such pair of functionals computes it.
The excluded relation therefore implies $A\in\core_X$. The reverse inclusion
holds because every image is in $\CD_X$. This proves
\eqref{eq:perfect-main}.

The unary requirements exclude every $S_{e,A}$ for $A\leT B_j$ outside the
core. The argument of Theorem~\ref{thm:partners} then proves
\eqref{eq:prescribed-main}. The singleton requirements ensure
$P\cap\core_X=\varnothing$. If distinct $D,E\in P$ satisfied $D\leT E$, the
set $D$ itself would lie in their common lower cone, hence in $\core_X$,
contradicting this exclusion. This proves pairwise incomparability.

Finally, the identity map from $(\CD_X,d)$ into ordinary Cantor space is
continuous by Lemma~\ref{lem:prefix}. Composing it with $F$ yields a continuous
injection from compact Cantor space into a Hausdorff space. It is therefore a
homeomorphism onto its image. The image is compact, hence closed, and has no
isolated points. Thus $P$ is perfect in the usual sense as well.
\end{proof}

\begin{remark}[A useful safeguard]
A perfect family of Turing minimal pairs does not alone disprove a least-degree
claim about its surrounding coarse class. The computable coarse class also
contains many noncomputable representatives. The later application separately
proves that its coarse class contains no computable representative.
\end{remark}

\section{Computable approximability forces a small core}
\label{sec:approximability}
We next produce a core to which the exact-pair theorem can be applied. The
following proof is independent of the older cone-avoiding compactness proof.
Its unrelativized conclusion is the result attributed to Igusa in
\cite[Theorem~4.3]{HJKS}; no novelty is claimed for that conclusion.

\begin{definition}
For an oracle $Z$, let $\AP_Z(X)$ denote the property
\begin{equation}\label{eq:AP}
 \forall\varepsilon>0\ \exists C\leT Z\quad \delta(C,X)<\varepsilon.
\end{equation}
The sequence of approximating sets need not be uniformly computable in $Z$.
For the computable oracle, this is the usual assertion that the coarse
computability bound is $1$.
\end{definition}

\begin{lemma}[Patching an approximate center]\label{lem:approxpatch}
If $\delta(C,Y)<\eta$, there is a finite modification $C_0$ of $C$ such that
$d(C_0,Y)<\eta$.
\end{lemma}
\begin{proof}
Choose $\eta'$ strictly between $\delta(C,Y)$ and $\eta$. For some $N$, all
prefixes of length $n>N$ have discrepancy less than $\eta'$. Put
$C_0=(Y\restr N)\patch C$. Prefixes of length at most $N$ have no errors,
and longer prefixes have no more errors than $C\triangle Y$. Consequently
$d(C_0,Y)\le\eta'<\eta$. The finite patch does not change the Turing degree.
\end{proof}

\begin{theorem}[Relative approximability bound]\label{thm:APcore}
If $\AP_Z(X)$, then $\core_X\subseteq\low Z$. In particular,
$\AP_{\emptyset}(X)$ implies $\core_X=\Comp$.
\end{theorem}
\begin{proof}
Fix $A\in\core_X$. Every member of $\CD_X$ computes $A$, so
$\CD_X=\bigcup_e S_{e,A}$. By the Baire category theorem, at least one $S_{e,A}$
is not nowhere dense. Hence it is dense in some nonempty open set $U$.
Choose $Y\in U$ and a positive rational $r$ with $B_d(Y,r)\subseteq U$.

By \eqref{eq:AP}, choose $C\leT Z$ with $\delta(C,X)<r/16$. Since
$\delta(X,Y)=0$, also $\delta(C,Y)<r/16$. Lemma~\ref{lem:approxpatch} gives
a finite modification $C_0\leT Z$ such that
\begin{equation}\label{eq:external-center}
 d(C_0,Y)<r/8.
\end{equation}
Crucially, $C_0$ is \emph{not} asserted to belong to $\CD_X$.

To compute $A(n)$ from $C_0$, search finite strings $\sigma$ and convergent
simulations $\Phi_e^\sigma(n)\downarrow=v$ subject to
\begin{equation}\label{eq:APsearch}
 \frac{|\{j<k:\sigma(j)\ne C_0(j)\}|}{k}<r/4
 \qquad(1\le k\le|\sigma|).
\end{equation}
This is an ordinary $C_0$-computable search.

For termination, the dense set $S_{e,A}$ meets $B_d(Y,r/8)\subseteq U$.
Choose $D$ in that intersection. Then
$d(D,C_0)\le d(D,Y)+d(Y,C_0)<r/4$, so a finite prefix of $D$ witnessing
convergence on input $n$ satisfies \eqref{eq:APsearch}.

For correctness, take any accepted $\sigma$ of length $L$ and form
$W=\sigma\patch Y$, patching onto the genuine description $Y$, not onto
$C_0$. Then $W\in\CD_X$. For $1\le k\le L$, the prefix triangle inequality,
\eqref{eq:external-center}, and \eqref{eq:APsearch} give
\[
 \rho_k(W\triangle Y)
 \le \frac{|\{j<k:\sigma(j)\ne C_0(j)\}|}{k}
       +\rho_k(C_0\triangle Y)
 <r/4+r/8=3r/8.
\]
After $L$ there are no additional errors, so
$d(W,Y)<3r/8<r$, by the finite-prefix identity. Thus
$W\in U\cap[\sigma]$. This nonempty open set meets $S_{e,A}$, forcing the
accepted value $v$ to be $A(n)$.

We have computed $A$ from $C_0\leT Z$. Since $A$ was any member of the core,
$\core_X\subseteq\low Z$. For computable $Z$, the opposite inclusion of all
computable sets follows from Lemma~\ref{lem:ideal}.
\end{proof}

\begin{remark}[Why the external-center argument matters]
A computable set can approximate $X$ to every fixed positive error tolerance
without any computable set being a coarse description of $X$. It would be
incorrect to insert such an approximate set directly into a proof that assumes
membership in $\CD_X$. The two different patches in the proof above separate
what the oracle computation can use from what the topological correctness
argument can use.
\end{remark}

\section{An explicit complete c.e. counterexample}
\label{sec:diagonal}
We now construct $X$ rather than simply citing the existence of a random or
generic example. Fix an acceptable enumeration $(\varphi_e)$ of partial
computable functions $\N\to\N$. It contains all total binary-valued computable
functions. Let
\begin{equation}\label{eq:columns}
 R_e=\{n:\nu_2(n+1)=e\}
     =\{2^e(2k+1)-1:k\in\N\},
\end{equation}
where $\nu_2$ is the exponent of $2$ dividing a positive integer. These sets
partition $\N$. Their densities and exact tail counts are
\begin{equation}\label{eq:column-densities}
 \rho(R_e)=2^{-e-1},\qquad
 \bigcup_{e\ge k}R_e=\{n:2^k\mid n+1\},\qquad
 \left|\Bigl(\bigcup_{e\ge k}R_e\Bigr)\cap[0,n)\right|
     =\left\lfloor\frac n{2^k}\right\rfloor.
\end{equation}

\begin{definition}[Guarded diagonal set]\label{def:X}
For $n\in R_e$, put
\begin{equation}\label{eq:explicit-X}
 n\in X\quad\Longleftrightarrow\quad
 \left[\forall m\le n\ \varphi_e(m)\downarrow\in\{0,1\}\right]
 \ \land\ \varphi_e(n)=0.
\end{equation}
\end{definition}

The finite-prefix convergence guard is the point of the construction. Without
it, a partial program might leave a noncomputable pattern inside one
positive-density column. With the guard, every non-total or nonbinary program
contributes only finitely many points to its column.

\begin{theorem}[Properties of the guarded diagonal set]\label{thm:X}
The set $X$ is computably enumerable and many-one complete. It has
$\AP_{\emptyset}(X)$, but it has no computable coarse description. Hence
\[
 X\eqT\emptyset',\qquad \core_X=\Comp,
 \qquad X\text{ is not coarsely computable}.
\]
\end{theorem}
\begin{proof}
\emph{Computable enumerability.} At stage $s$, examine each $n\le s$, compute
its unique column index $e$, and simulate the finitely many computations
$\varphi_e(0),\ldots,\varphi_e(n)$ for $s$ steps. Enumerate $n$ if all have
halted with binary values and the last has value zero. An enumerated witness
remains valid forever. Conversely, every witness to \eqref{eq:explicit-X}
is discovered at a finite stage. This is a c.e. enumeration of $X$.

\emph{Individual columns are computable.} Fix $e$. If $\varphi_e$ is total
and binary-valued, then on $R_e$,
\[
 X(n)=1-\varphi_e(n).
\]
In this case $X\cap R_e$ is computable. Otherwise let $m_e$ be the least input
on which $\varphi_e$ is undefined or has a nonbinary value. For every
$n\ge m_e$, the guard in \eqref{eq:explicit-X} fails, so
$X\cap R_e\subseteq[0,m_e)$. This column is finite and hence computable.
These alternatives prove computability for each fixed column, without an
effective procedure deciding which alternative holds.

\emph{Arbitrarily accurate computable approximations.} For fixed $k$, let
\[
 C_k=X\cap\bigcup_{e<k}R_e.
\]
It is a finite union of computable sets, so it is computable. All errors of
$C_k$ lie in the tail union $\bigcup_{e\ge k}R_e$. Thus
\begin{equation}\label{eq:approx-rate}
 \delta(C_k,X)\le2^{-k}.
\end{equation}
For any positive $\varepsilon$, choose $k$ with $2^{-k}<\varepsilon$.
This proves $\AP_{\emptyset}(X)$. It does not assert that indices for the
$C_k$ can be obtained effectively from $k$.

\emph{No computable coarse description.} If $C$ is a computable set, its
characteristic function is $\varphi_e$ for some total binary index $e$.
Then $X(n)\ne C(n)$ for every $n\in R_e$. Therefore
\[
 \liminf_{n\to\infty}\rho_n(X\triangle C)
 \ge \rho(R_e)=2^{-e-1}>0.
\]
So $C$ is not a coarse description. A natural-number-valued computable coarse
description would yield a binary one by the normalization in
Section~\ref{sec:definitions}, and is excluded as well.

\emph{Many-one completeness.} Let
$K=\{p:\varphi_p(p)\downarrow\}$ be the usual halting set. By the parameter
 theorem, obtain computably from $p$ an index $e(p)$ of the following program:
on any input, first wait for $\varphi_p(p)$ to halt, and then return zero.
If $p\in K$, this program is total and constantly zero. If $p\notin K$, it is
nowhere defined. Put $n(p)=2^{e(p)}-1$, which belongs to $R_{e(p)}$. Then
\[
 p\in K\quad\Longleftrightarrow\quad n(p)\in X.
\]
The map $p\mapsto n(p)$ is total computable, giving $K\le_m X$. As $X$ is
c.e., $X\leT K$, so $X\eqT\emptyset'$.

Finally, Theorem~\ref{thm:APcore} applied to the proved approximability gives
$\core_X=\Comp$.
\end{proof}

\begin{remark}[An index-level subtlety]
The displayed formula for $C_k$ does not give a uniform algorithm for its
characteristic function from $k$. The proof of each column's computability
uses either a genuinely total program or the existence of a fixed finite set.
Recognizing the appropriate case uniformly would decide information about
program totality. No such recognition is used in this construction or in
Theorem~\ref{thm:APcore}.
\end{remark}

\section{Minimal pairs, no least degree, and uncountable coinitiality}
\label{sec:consequences}
From now on $X$ denotes the explicit set of Definition~\ref{def:X}.

\begin{lemma}[The coarse classes are nonzero]\label{lem:nonzero}
Neither $[X]_{\nc}$ nor $[X]_{\uc}$ contains a computable representative,
including a computable natural-number-valued function.
\end{lemma}
\begin{proof}
If a computable $g$ belonged to either class, then $X\le_{\nc}g$. Since $g$
is its own computable coarse description, the reduction would give a
computable coarse description of $X$, contrary to Theorem~\ref{thm:X}.
\end{proof}

\begin{theorem}[The requested minimal-pair certificate]\label{thm:minpair}
There is $D\equiv_{\uc}X$ such that $\degT D$ and $\mathbf0'$ form a Turing
minimal pair. In particular both $[X]_{\nc}$ and $[X]_{\uc}$ have no least
Turing degree.
\end{theorem}
\begin{proof}
Theorem~\ref{thm:partners} gives $D\in\CD_X$ with
\[
 (\low D)\cap(\low X)=\core_X=\Comp.
\]
Lemma~\ref{lem:sameclass} gives $D\equiv_{\uc}X$, and
Lemma~\ref{lem:nonzero} makes $D$ noncomputable. The degree of $X$ is
$\mathbf0'$. Thus the two degrees are nonzero and their only common lower
degree is $\mathbf0$, exactly the minimal-pair condition.

If $g$ were a least representative of either coarse class, it would be
computable from both $X$ and $D$. Apply the common-lower-cone equality to the
graph of $g$ when $g$ is a function. Then $g$ would be computable, contradicting
Lemma~\ref{lem:nonzero}.
\end{proof}

\begin{corollary}[Perfect strengthening]\label{cor:perfectminimal}
There is a perfect set $P\subseteq\CD_X$ whose members form Turing minimal
pairs with each other and, individually, with $\emptyset'$. Every member is
uniformly coarsely equivalent to $X$. In particular, every member is Turing
incomparable with $\emptyset'$ and is not $\Delta^0_2$.
\end{corollary}
\begin{proof}
Apply Theorem~\ref{thm:perfect} with prescribed oracle $B_0=X$. The core is
$\Comp$, and no member of $\CD_X$ is computable. The exact-pair equalities
therefore become minimal-pair equalities. A nonzero degree forming a minimal
pair with $\mathbf0'$ can neither be below nor above $\mathbf0'$; in
particular no representative is computable from $\emptyset'$.
\end{proof}

\begin{definition}
For $r\in\{\nc,\uc\}$, write
\[
 \Spec([X]_r)=\{\degT g:g\equiv_r X\}.
\]
A subset $Q$ of this spectrum is \emph{coinitial} if every degree in the
spectrum bounds some degree in $Q$. This is the order-theoretic notion of
coinitiality, not a claim about a computable listing.
\end{definition}

\begin{theorem}[No countable coinitial family]\label{thm:coinitial}
For either $r\in\{\nc,\uc\}$ and every countable sequence
$g_j\equiv_r X$, there is $D\in\CD_X$ forming a Turing minimal pair with
every $g_j$. Consequently $\Spec([X]_r)$ has no countable coinitial subset.
Indeed one may choose a perfect family of such $D$, pairwise forming minimal
pairs with each other.
\end{theorem}
\begin{proof}
Let $B_j$ be a set coding the graph of $g_j$, so that $B_j\eqT g_j$.
Theorem~\ref{thm:partners}, or its perfect strengthening, gives $D\in\CD_X$
with
\[
 (\low D)\cap(\low {B_j})
   =\core_X\cap(\low {B_j})=\Comp
 \quad\text{for every }j.
\]
All $g_j$ and $D$ are noncomputable by Lemma~\ref{lem:nonzero}. Therefore the
equalities are minimal-pair equalities. In particular $g_j\nleT D$ for every
$j$. Since $D\equiv_r X$, no degree in the proposed countable family lies
below this member of the spectrum. The family is not coinitial.
\end{proof}

\begin{remark}[Why this is stronger than the original target]
A least degree would give a one-element coinitial family. The theorem rules
out even a countably infinite one, and does so with an oracle Turing
incomparable with each proposed member. It does not exclude an uncountable
set of minimal degrees. No exact cardinal value for the coinitiality beyond
uncountability is asserted here.
\end{remark}

\subsection{The representative spectrum is upward closed}
For completeness, the familiar sparse-coding observation in the source
report can be checked directly. Suppose $Y\equiv_r X$ is binary and
$Y\leT B$. Define $W$ by putting $W(2^n)=B(n)$ and keeping $W(k)=Y(k)$ at
all other positions. The powers of two have density zero, so
$W\equiv_{\uc}Y$. Also $B\leT W$ by reading those positions, while
$W\leT B$ because $Y\leT B$. Thus $\degT W=\degT B$.

The same argument works for function representatives by first choosing a
set $B$ of the desired higher degree and storing its bits on the powers of
two. This shows that both full representative spectra are upward closed.
There is no conflict with the preceding theorems: an upward-closed family of
Turing degrees need not have a least member or a countable coinitial basis.

\section{Which nonuniform coarse classes have least representatives?}
\label{sec:classification}
The least-degree property has a short structural characterization. We give
it to clarify exactly why the earlier nonsurjectivity results already bear
on C1. The proof uses a simple robust coding, with no uniform decoder claimed.

For $A\subseteq\N$, let $J(A)(0)=0$ and define
\begin{equation}\label{eq:J}
 J(A)(k)=A(n)\qquad\text{when }2^n\le k<2^{n+1}.
\end{equation}
Thus bit $A(n)$ is repeated on a block of length $2^n$.

\begin{lemma}[Nonuniform robust recovery]\label{lem:J}
$J(A)\eqT A$, and every coarse description of $J(A)$ computes $A$.
Consequently
\[
 \core_{J(A)}=\low A.
\]
\end{lemma}
\begin{proof}
Computing $J(A)$ from $A$ uses the block index. Conversely
$A(n)=J(A)(2^n)$, proving Turing equivalence.

Let $D$ be a binary coarse description of $J(A)$. Compute $b(n)$ as the
majority value of $D$ on $[2^n,2^{n+1})$, resolving ties by zero. For all large
$n$,
\[
 |(D\triangle J(A))\cap[0,2^{n+1})|
       <\tfrac14\,2^{n+1}=2^{n-1}.
\]
This is fewer than half the positions in the $n$th block, so $b(n)=A(n)$ for
all sufficiently large $n$. Thus $A$ is a finite modification of the
$D$-computable set $b$, and $A\leT D$. The finite correction depends on $D$;
there is no asserted uniform recovery of its cutoff or values.

All sets below $A$ are therefore in the core. The reverse inclusion follows
from $J(A)\in\CD_{J(A)}$ and $J(A)\eqT A$.
\end{proof}

\begin{theorem}[Least-representative characterization]\label{thm:leastclass}
A nonuniform coarse-equivalence class of total functions has a least Turing
degree if and only if it is $[J(A)]_{\nc}$ for some set $A$. In that case its
least degree is $\degT A$.
\end{theorem}
\begin{proof}
First suppose $g\equiv_{\nc}J(A)$. Since $J(A)\le_{\nc}g$, the oracle $g$
computes a coarse description $D$ of $J(A)$. Lemma~\ref{lem:J} yields
$A\leT D\leT g$. The representative $J(A)$ itself has degree $\degT A$, so
that degree is least in the class.

Conversely suppose the class has a least representative $g$, and choose a
set $A\eqT g$ by graph coding. Every coarse description $h$ of $g$ is itself
uniformly coarsely equivalent to $g$, by the disagreement triangle inequality.
Leastness therefore gives $g\leT h$, and hence $A\leT h$. Thus every coarse
description of $g$ computes $J(A)$ itself, proving $J(A)\le_{\nc}g$.

In the other direction every coarse description of $J(A)$ computes $A$ by
Lemma~\ref{lem:J}, and hence computes $g$ itself. Therefore
$g\le_{\nc}J(A)$. This proves equality of the classes and identifies the least
degree.
\end{proof}

The map $\degT A\mapsto[J(A)]_{\nc}$ is an order embedding: if $A\leT B$,
every description of $J(B)$ computes $B$ and then $J(A)$; if
$J(A)\le_{\nc}J(B)$, apply the reduction to $J(B)$ itself to get
$A\leT J(B)\eqT B$. Thus the characterization identifies exactly its image.

\begin{remark}[Uniformity boundary]
The theorem characterizes \emph{nonuniform} coarse classes. The eventual
majority decoding above does not provide one exact decoder for $A$ that works
on every description, and the converse proof does not establish a uniform
coarse equivalence. No corresponding image characterization for uniform
coarse classes is asserted. The negative uniform-class results in
Section~\ref{sec:consequences} do not depend on such a characterization.
\end{remark}

\section{Relativized examples with any principal core}
\label{sec:relative}
The construction is not restricted to common information of computable
degree. Fix any set $Z$. Using an acceptable enumeration of partial
$Z$-computable functions, define $X_Z$ by the guarded formula
\eqref{eq:explicit-X}, with $\varphi_e^Z$ in place of $\varphi_e$. The proof of
Theorem~\ref{thm:X} relativizes line by line:
\begin{equation}\label{eq:relativeX}
 X_Z\text{ is c.e. in }Z,\qquad X_Z\eqT Z',\qquad
 \AP_Z(X_Z),
\end{equation}
and $X_Z$ has no $Z$-computable coarse description. The reductions proving
$X_Z\eqT Z'$ have the same form: $Z'$ enumerates membership in any $Z$-c.e.
set, and the parameterized all-or-nothing programs reduce $Z'$ to $X_Z$.

Let
\begin{equation}\label{eq:Y}
 Y_Z=J(Z)\oplus X_Z,
 \quad Y_Z(2n)=J(Z)(n),\quad Y_Z(2n+1)=X_Z(n).
\end{equation}

\begin{theorem}[Arbitrary principal cores]\label{thm:relative}
For every $Z$ the set $Y_Z$ has
\[
 Y_Z\eqT Z',\qquad \core_{Y_Z}=\low Z,
 \qquad Y_Z\text{ has no }Z\text{-computable coarse description}.
\]
Every representative of $[Y_Z]_{\nc}$ or $[Y_Z]_{\uc}$ has Turing degree
strictly above $\degT Z$. Both spectra have no countable coinitial family.
There is a perfect family of representatives whose distinct pairs have
greatest common lower degree exactly $\degT Z$.
\end{theorem}
\begin{proof}
Taking the odd half of $Y_Z$ computes $X_Z$, and both halves are computable
from $Z'$, proving $Y_Z\eqT Z'$. If $D$ is a coarse description of $Y_Z$,
its even half is a coarse description of $J(Z)$: up to $n$ positions, its
errors are bounded by the total errors of $D$ up to $2n$, so their proportion
is at most twice that total proportion. Lemma~\ref{lem:J} gives $Z\leT D$.
Thus $\low Z\subseteq\core_{Y_Z}$.

Given $\varepsilon>0$, choose a $Z$-computable $C$ with
$\delta(C,X_Z)<\varepsilon$. The join $J(Z)\oplus C$ is $Z$-computable, and
only its odd half can disagree with $Y_Z$. Its upper discrepancy is
$\tfrac12\delta(C,X_Z)<\varepsilon/2$. Therefore $\AP_Z(Y_Z)$ holds.
Theorem~\ref{thm:APcore} gives the opposite inclusion of cores, proving
$\core_{Y_Z}=\low Z$.

A $Z$-computable coarse description of $Y_Z$ would have a $Z$-computable odd
half describing $X_Z$, which is impossible. If $g$ belongs to either coarse
class, $Y_Z\le_{\nc}g$, so $g$ computes a coarse description of $Y_Z$, and
hence computes $Z$. If also $g\leT Z$, that same description would be
$Z$-computable, a contradiction. Thus every such $g$ is strictly above $Z$.

For any countable family of representatives $g_j$, code them by oracles
$B_j$. Theorem~\ref{thm:partners} gives a description $D$ whose common lower
cone with each $B_j$ is $\low Z$, since every $B_j$ computes $Z$. Therefore
$g_j\nleT D$: otherwise $g_j$ would be computable from $Z$, contrary to what
was just proved. This excludes countable coinitiality. Theorem~\ref{thm:perfect}
gives the perfect strengthening, with common lower cones exactly $\low Z$.
\end{proof}

This result concerns ordinary Turing degrees above $\degT Z$; it does not
change the underlying definition of coarse reducibility. The parameter $Z$
is used to construct an example whose common information is exactly that
oracle's lower cone.

\section{Proof audit, limitations, and continuation targets}
\label{sec:audit}
\subsection{Dependency ledger}
The logical dependencies are short enough to expose explicitly:
\[
\begin{gathered}
 \text{prefix-density inequalities}
 \Longrightarrow \text{complete description space}
 \Longrightarrow \text{local recovery},\\
 \text{local recovery}+\text{Baire category}
 \Longrightarrow \text{simultaneous exact partners},\\
 \text{local recovery}+\text{finite-level fusion}
 \Longrightarrow \text{perfect exact pairs},\\
 \text{approximability}+\text{Baire category}+\text{prefix search}
 \Longrightarrow \text{small core},\\
 \text{guarded c.e. diagonalization}+\text{small core}
 \Longrightarrow \text{the explicit negative examples}.
\end{gathered}
\]
The many-one completeness statement additionally uses the standard parameter
 theorem for an acceptable enumeration. No randomness basis theorem, jump
inversion theorem, classical priority minimal-pair construction, or
uncountable effective listing is assumed.

\subsection{Where incorrect shortcuts were excluded}
The density metric $\delta$ cannot support the category proof: it vanishes
throughout one description class. The supremum-prefix metric $d$ is used
instead, and its completeness and separability are proved.

Individual cone avoidance is not silently upgraded to simultaneous avoidance.
The nowhere-density calculation precedes the countable Baire argument. The
perfect construction explicitly handles all ordered pairs of distinct nodes
at each level; separate pairwise existence assertions would not produce a
single compatible perfect family.

The set used as a computable approximate center in
Theorem~\ref{thm:APcore} need not be a description. The algorithm queries that
center, but the correctness proof patches onto a genuine description. This
separation is indispensable.

The individual column programs in the c.e. construction are not presented as
a uniformly computable sequence. The finite-correction argument in robust
block decoding is also explicitly nonuniform. Neither step recovers unknown
indices or stabilization stages effectively.

Finally, membership in the same coarse class is proved before applying the
minimal-pair criterion, and nonzero coarse degree is proved independently.
Ordinary noncomputability of the chosen binary set would not suffice.

\subsection{Effectivity and formal verification}
The set $X$ is given by an explicit c.e. enumeration, with a proof that it is
many-one complete. The additional descriptions and the perfect set are
obtained by classical countable choices and Baire/fusion arguments. We do not
provide a computable procedure for the perfect-set code, an oracle bound on
that code, or a low arithmetical bound on its paths. In fact the selected paths
forming minimal pairs with $\emptyset'$ cannot be $\Delta^0_2$.

No Lean executable was available in the working environment, and no Lean
proof is supplied or claimed to compile. The source report is not treated as
an axiom library. The mathematical proofs in this note are conventional
proofs; the program in the archive checks only explicitly stated finite
lemmas. Those checks are not evidence that the infinite theorems have been
kernel verified.

\subsection{The next mathematical questions}
A meaningful continuation would determine how effectively the perfect family
can be selected from a presentation of $X$ and its core. The proof exposes
the relevant obstacles: deciding that information lies outside the core,
locating open rectangles outside closures of bad computation relations, and
making compatible shrinkings. None of these choices has been assigned an
arithmetical bound here. This is a newly proposed refinement, not a claim that
it is a verified literature-open problem.

The nonprincipal-core case also deserves investigation. The perfect theorem
already realizes any core arising from a coarse-description class as a
common lower-cone ideal of pairs inside that same class. What additional
restrictions characterize the ideals that arise this way is not determined
here.

The effective-dense analogue cannot be inferred. An effective-dense
description must either answer correctly or explicitly omit an answer; it
cannot silently make wrong assertions. Density-zero modifications preserve
coarse descriptions but do not in general preserve effective-dense
descriptions. Thus the argument leaves C2 unchanged.

\appendix
\section{Finite checks and reproduction}
\label{app:checks}
The archive contains one standard-library Python program,
\path{checks/check_finite_lemmas.py}. It uses exact integers and rational
fractions, with no floating-point density comparisons. The run packaged with
this note completed \textbf{302,579 finite checks}, all passing.

\begin{center}
\begin{tabular}{@{}lr@{}}
\toprule
Finite statement or model & Cases\\
\midrule
Metric axioms on binary strings of length at most five & 37,449\\
Exact finite-prefix patch identity & 43,435\\
External-center patch inequality & 9,021\\
Positive-density column counts & 53,261\\
Tail-column counts & 53,261\\
Column truncation error bounds & 12,282\\
Guarded diagonalization on finite partial tables & 13,120\\
Monotonicity of finite delayed enumerations & 27,884\\
Strict-majority block recovery & 52,866\\
\midrule
Total & 302,579\\
\bottomrule
\end{tabular}
\end{center}

The finite partial tables use a distinguished value to mean \emph{permanently
undefined within the model}, not a claim that a real program has failed to
halt. The tests therefore do not decide totality. Similarly, finite column
counts support the displayed arithmetic identities, not the Turing
computability of individual infinite columns.

To rerun the checks and compile the article from the archive root:
\begin{lstlisting}
python checks/check_finite_lemmas.py --output checks/results.json
pdflatex -interaction=nonstopmode -halt-on-error coarse_degree_attack.tex
pdflatex -interaction=nonstopmode -halt-on-error coarse_degree_attack.tex
\end{lstlisting}
Python 3.9 or later and a standard TeX Live installation suffice. The provided
\path{build.sh} also performs these steps in an isolated build directory.
The shipped PDF was compiled locally and rendered for visual inspection.
The source report is included unchanged under \path{source/} for provenance;
it is not needed to compile this article.

\begin{thebibliography}{9}
\addcontentsline{toc}{section}{References}
\bibitem{Survey}
\emph{Turing Degrees: A Unified Research Report}.
Supplied file \path{turing_degrees_unified.tex}, unified edition
18 September 2026. In particular entry C1 and its minimal-pair diagnostic.
An unchanged copy is included in this archive under \path{source/}.

\bibitem{Gerdes}
Peter M. Gerdes.
\emph{Comparing Notions of Dense Computability on $\omega^\omega$ and
$2^\omega$}.
arXiv:2508.06925v1, 9 August 2025. Definitions 2.3--2.4 and Question 7.
\url{https://arxiv.org/abs/2508.06925v1}.

\bibitem{HJKS}
Denis R. Hirschfeldt, Carl G. Jockusch, Jr., Rutger Kuyper, and Paul E. Schupp.
\emph{Coarse Reducibility and Algorithmic Randomness}.
Journal of Symbolic Logic \textbf{81}(3) (2016), 1028--1046.
Preprint arXiv:1505.01707. In particular Definition 3.1 and Theorem 4.3.
\url{https://arxiv.org/abs/1505.01707}.
Author's publication record:
\url{https://www.math.uchicago.edu/~drh/Papers/coarsereducibility.html}.
\end{thebibliography}
\end{document}
